%==============================================================================
% Chapter 6 — Deep Learning Fundamentals
% Curriculum: Phase 2, Chapter 6 of docs/AI_ML_Master_Checklist.md
% Companion code: projects/06_deep_learning/
%
% Section skeleton only. Figures for this chapter go in theory/figures/.
%==============================================================================
\chapter{Deep Learning Fundamentals}
\label{ch:deep-learning}

The pivot of the book.\projectref{projects/06\_deep\_learning/} Everything in \cref{ch:classical-ml} either had a closed
form or a bespoke fitting procedure; from here on there is one procedure ---
write down a differentiable loss and descend its gradient --- and the modelling
work moves into choosing the architecture.

The chapter closes on two things that Phase~2 cannot proceed without. Sequence
models --- recurrence, gating, and the encoder/decoder bottleneck --- are what
make the Transformer of \cref{sec:transformer} read as a solution rather than a
declaration. High-performance computing is what makes any of
\cref{ch:gen-ai} runnable at all: the constraint on a modern model is GPU
memory, and that is arithmetic you can do in advance rather than discover
through a crash.

The mathematics is already in place. Backpropagation is \cref{eq:chain-rule}
evaluated in reverse topological order, and the optimisers are variations on
descending the gradient of \cref{eq:taylor}. The goal of the chapter is to
write both out by hand for a multilayer perceptron, implement them from scratch
against a finite-difference check, and only then reach for a framework.

%==============================================================================
\section{The Multilayer Perceptron}
\label{sec:mlp}

\subsection{Nodes, Weights and Biases}

\subsection{The Forward Pass}

\subsection{Backpropagation by Hand}

\subsection{Universal Approximation and Its Limits}

%==============================================================================
\section{Activation Functions}
\label{sec:activations}

\subsection{Sigmoid and Tanh}

\subsection{ReLU and Its Variants}

\subsection{GELU}

\subsection{Softmax}

%==============================================================================
\section{Loss Functions}
\label{sec:losses}

\subsection{Mean Squared Error}

\subsection{Binary and Categorical Cross-Entropy}

\subsection{Hinge Loss}

%==============================================================================
\section{Optimization Algorithms}
\label{sec:optimizers}

\subsection{Stochastic Gradient Descent}

\subsection{Momentum}

\subsection{RMSProp}

\subsection{Adam and the Bias-Corrected Moment Estimates}

\subsection{Learning Rate Schedules}

%==============================================================================
\section{Regularization}
\label{sec:regularization}

\subsection{Dropout}

\subsection{Batch Normalization}

\subsection{Early Stopping}

\subsection{Weight Decay and Its Relation to Ridge}

%==============================================================================
\section{Training in Practice}
\label{sec:training-practice}

\subsection{Weight Initialisation}

\subsection{Vanishing and Exploding Gradients}

\subsection{Diagnosing a Training Run}

%==============================================================================
\section{Sequence Models}
\label{sec:sequence-models}

% The bridge to \cref{ch:gen-ai}. Attention is far harder to motivate if the
% fixed-context bottleneck it removes was never felt, so this section exists to
% make the Transformer feel inevitable rather than arbitrary.

\subsection{Recurrent Neural Networks}

\subsection{Backpropagation Through Time}

\subsection{Vanishing Gradients Over Time}

% The same phenomenon as \cref{sec:training-practice} but compounded by weight
% sharing across timesteps: the gradient carries a factor of the recurrent
% weight matrix raised to the sequence length.

\subsection{LSTM Gating}

\subsection{Gated Recurrent Units}

\subsection{Sequence-to-Sequence Models}

\subsection{The Fixed-Context Bottleneck}

\subsection{Additive Attention}

% Bahdanau attention: let the decoder form a weighted sum over all encoder
% states instead of a single final vector. Removing the recurrence from this
% idea is exactly \cref{sec:transformer}.

%==============================================================================
\section{Efficient Training and High-Performance Computing}
\label{sec:hpc}

% What separates "train an MLP on a laptop" from "fine-tune a language model",
% and therefore a prerequisite for \cref{ch:gen-ai}. The binding constraint is
% almost always memory, and it is predictable with arithmetic rather than by
% trial and error.

\subsection{The GPU Execution Model}

% SIMT, warps, occupancy and the memory hierarchy --- enough to explain why
% small batches leave the device idle and why data loading is so often the real
% bottleneck.

\subsection{Memory Arithmetic}

% Parameters + gradients + optimizer states + activations. Adam holds two
% moments per parameter, so an FP32 model costs roughly 16 bytes per parameter
% before activations. Work out whether a model fits before launching the job.

\subsection{Mixed Precision and Loss Scaling}

% FP16 versus BF16, why FP16 needs loss scaling, and how this connects back to
% \cref{sec:numerical-stability}.

\subsection{Gradient Accumulation}

\subsection{Activation Checkpointing}

\subsection{Data Parallelism: DDP}

\subsection{Sharded Training: ZeRO and FSDP}

\subsection{Tensor, Pipeline and Model Parallelism}

\subsection{Profiling and Finding the Real Bottleneck}

\subsection{Cluster Scheduling and Checkpoint Discipline}

% SLURM job submission, multi-GPU and multi-node launches, and writing
% checkpoints such that a pre-empted job resumes rather than restarts.
